\documentclass[10pt,letterpaper,sans]{moderncv}
\moderncvstyle{banking}
\moderncvcolor{blue}

\usepackage[scale=0.81]{geometry}
\nopagenumbers{}
\setlength{\hintscolumnwidth}{2.25cm}

\name{Krishnaa}{Vadivel}
\title{GPU Software Engineer | High-Performance Computing | Scientific ML}
\email{srikrishnaa.careers@gmail.com}
\phone[mobile]{516-853-5663}
\social[linkedin]{sri-krishnaa-ece-phd}
\extrainfo{\href{https://krishnaa423.github.io/personal-website/}{Website} \textbar{} \href{https://github.com/krishnaa423}{GitHub: krishnaa423} \textbar{} \href{https://leetcode.com/u/krishnaa42342}{LeetCode: krishnaa42342} \textbar{} \href{https://hub.docker.com/u/krishnaa42342}{Docker Hub: krishnaa42342}}

\begin{document}
\makecvtitle

\section{Summary}
\cvitem{}{GPU-focused computational scientist and software engineer building accelerator-enabled applications in Python and C++ for large-scale simulation and data-processing workloads. Experience spans CUDA, ROCm, MPI, OpenMP, national-lab supercomputing systems, and machine-learning models that accelerate scientific computing pipelines.}

\section{Experience}
\cventry{2021--present}{PhD Research Assistant}{Yale University, Electrical Engineering}{}{}{\begin{itemize}%
\item Developed distributed scientific software in Python and C++ for first-principles simulation workloads as part of an INCITE award using exascale resources, with MPI, OpenMP, CUDA, ROCm, PETSc, and SLEPc on Frontier at OLCF, Perlmutter at LBNL, and related national-lab clusters.
\item Ported and optimized scientific-computing workloads for accelerator-enabled HPC environments, translating many-body and excited-state physics methods into scalable GPU-backed implementations.
\item Built PyTorch and PyTorch Geometric models, including equivariant graph neural networks, to accelerate self-trapped-exciton calculations and related materials simulations.
\item Packaged scientific software in Docker containers and used tools such as TotalView and Linaro DDT to debug distributed applications across heterogeneous compute environments.
\item Co-authored a 2025 \textit{Journal of the American Chemical Society} paper and presented APS research in 2024, 2025, and 2026, grounding GPU and HPC work in production research workflows.
\end{itemize}}
\cventry{2019--2021}{Software and Embedded Systems Engineer}{Probe Technologies Inc. / Weatherford Wireline Services}{}{}{\begin{itemize}%
\item Developed CUDA-based software for high-volume acoustic waveform processing, visualization, and field-tool diagnostics in geological well logging systems.
\item Built full-stack server applications with C\#, ASP.NET, and Microsoft SQL Server for database design, real-time remote tool testing, and calibration of temperature and pressure sensors.
\item Applied finite-element analysis with dolfinx for frequency analysis of steel pipes in well environments to support physics-driven engineering studies.
\end{itemize}}
\cventry{2018--2019}{Photonics Technical Writer}{FindLight}{}{}{\begin{itemize}%
\item Wrote clear technical content on lasers, optics, and photonics devices for engineering audiences and product-focused readers.
\end{itemize}}

\section{Selected Projects}
\cvitem{Accelerator-Enabled Exciton-Phonon Simulation}{Developed distributed first-principles calculations for excited-state and exciton-phonon problems with CUDA, ROCm, MPI, and OpenMP on leadership-class GPU clusters.}
\cvitem{Scientific ML on GPU Systems}{Developed PyTorch Geometric equivariant models to accelerate molecular and excitonic calculations in GPU-backed scientific workflows.}
\cvitem{CUDA Containers for HPC Software}{Packaged scientific applications in CUDA-ready Docker environments for portable deployment across local and HPC systems.}
\cvitem{Matrix-Multiply AI Accelerator}{Built a Verilog-based matrix-multiply accelerator project for AI-oriented datapath design and hardware-aware compute fundamentals.}

\section{Education}
\cventry{2021--present}{PhD, Electrical Engineering}{Yale University}{}{}{}
\cventry{2023}{MPhil, Electrical Engineering}{Yale University}{}{}{}
\cventry{2023}{MS, Electrical Engineering}{Yale University}{}{}{}
\cventry{2017--2018}{MEng, Electrical and Computer Engineering}{Cornell University}{}{}{}
\cventry{2013--2017}{BS, Electrical and Computer Engineering}{Cornell University}{}{}{}

\section{Technical Skills}
\cvitem{Programming}{Python, C, C++, CUDA, JavaScript, Bash, SQL, C\#}
\cvitem{GPU / HPC}{CUDA, ROCm, MPI, OpenMP, PETSc, SLEPc, TotalView, Linaro DDT, Frontier, Perlmutter}
\cvitem{Machine Learning}{PyTorch, PyTorch Geometric, scikit-learn, equivariant graph neural networks}
\cvitem{Systems}{Linux, macOS, Windows, Docker, Docker Compose, Kubernetes, gcloud, CMake}
\cvitem{Domains}{GPU computing, high-performance computing, first-principles materials modeling, many-body theory, scientific software}

\end{document}
